%%%%%%%%%%%%%%%%
%%% exod 6 %%%
%%%%%%%%%%%%%%%%

\Chapter{6}

\Verse{1}
{
	\hRJE{וַיֹּאמֶר יְהֹוָה אֶל מֹשֶׁה עַתָּה תִרְאֶה אֲשֶׁר אֶעֱשֶׂה לְפַרְעֹה כִּי בְיָד חֲזָקָה יְשַׁלְּחֵם וּבְיָד חֲזָקָה יְגָרְשֵׁם מֵאַרְצוֹ׃}
}
{
	\eRJE{
		And YHWH said to Moses, “Now you’ll see what I shall do
		to Pharaoh, because with a strong hand he’ll let them go,
		and with a strong hand he’ll drive them from his land.”
	}
}
{
%	\footnote{
%		with a strong hand he’ll let them go. But with whose strong hand? God’s
%		or Moses’ own? The probable meaning is: God’s. But the ambiguity is rich
%		because either meaning is instructive.
%	}
}

\Verse{2}
{
	\hP{וַיְדַבֵּר אֱלֹהִים אֶל מֹשֶׁה וַיֹּאמֶר אֵלָיו אֲנִי יְהֹוָה׃}
}
{
	\eP{
		AND I APPEARED And {\God} spoke to Moses and said to him, “I
		am YHWH.
	}
}
{
}

\Verse{3}
{
	\hP{וָאֵרָא אֶל אַבְרָהָם אֶל יִצְחָק וְאֶל יַעֲקֹב בְּאֵל שַׁדָּי וּשְׁמִי יְהֹוָה לֹא נוֹדַעְתִּי לָהֶם׃}
}
{
	\eP{
		And I appeared to Abraham, to Isaac, and to Jacob as El
		Shadday, and I was not known to them by my name, YHWH.
	}
}
{
%	\footnote{
%		I was not known to them by my name, YHWH. But God was known to them by
%		this name (Gen 13:4; 15:2; 22:14,16; 26:22,25; 28:13). I have discussed
%		the significance of the understanding of this problem in critical
%		scholarship elsewhere (Who Wrote the Bible?). It has to do with the case
%		for multiple sources from which the Torah was composed. But in terms of
%		the final composition, the full Torah as we know it, how shall we
%		understand this verse? Perhaps covenantally: In the Noahic covenant (Gen
%		9:8–9), God is called ’lhîm. In the Abrahamic covenant, God is known as
%		El Shadday (Gen 17:1–2). But the covenant that will be made with this
%		generation at Sinai will be with the use of the name YHWH (Exod 20:2).
%		In support of this understanding: the very next words here are “And I
%		also established my covenant with them.”
%	}
}

\Verse{4}
{
	\hP{וְגַם הֲקִמֹתִי אֶת בְּרִיתִי אִתָּם לָתֵת לָהֶם אֶת אֶרֶץ כְּנָעַן אֵת אֶרֶץ מְגֻרֵיהֶם אֲשֶׁר גָּרוּ בָהּ׃}
}
{
	\eP{
		And I also established my covenant with them, to give them
		the land of Canaan, the land of their residences, in which
		they resided.
	}
}
{
}

\Verse{5}
{
	\hP{וְגַם אֲנִי שָׁמַעְתִּי אֶת נַאֲקַת בְּנֵי יִשְׂרָאֵל אֲשֶׁר מִצְרַיִם מַעֲבִדִים אֹתָם וָאֶזְכֹּר אֶת בְּרִיתִי׃}
}
{
	\eP{
		And also: I’ve heard the cry of the children of Israel as
		Egypt is enslaving them, and I’ve remembered my covenant.
	}
}
{
}

\Verse{6}
{
	\hP{לָכֵן אֱמֹר לִבְנֵי יִשְׂרָאֵל אֲנִי יְהֹוָה וְהוֹצֵאתִי אֶתְכֶם מִתַּחַת סִבְלֹת מִצְרַיִם וְהִצַּלְתִּי אֶתְכֶם מֵעֲבֹדָתָם וְגָאַלְתִּי אֶתְכֶם בִּזְרוֹעַ נְטוּיָה וּבִשְׁפָטִים גְּדֹלִים׃}
}
{
	\eP{
		Therefore, say to the children of Israel: ‘I am YHWH, and
		I shall bring you out from under Egypt’s burdens, and I
		shall rescue you from their toil, and I shall redeem you
		with an outstretched arm and with tremendous judgments,
	}
}
{
}

\Verse{7}
{
	\hP{וְלָקַחְתִּי אֶתְכֶם לִי לְעָם וְהָיִיתִי לָכֶם לֵאלֹהִים וִידַעְתֶּם כִּי אֲנִי יְהֹוָה אֱלֹהֵיכֶם הַמּוֹצִיא אֶתְכֶם מִתַּחַת סִבְלוֹת מִצְרָיִם׃}
}
{
	\eP{
		and I shall take you to me as a people, and I shall become
		your {\God}, and you’ll know that I am YHWH, your {\God}, who is
		bringing you out from under Egypt’s burdens.
	}
}
{
%	\footnote{
%		you’ll know that I am YHWH, your God, who is bringing you out. The
%		dominant theme of Exodus, YHWH’s becoming known, is connected explicitly
%		here to the rescue from slavery and to the “judgments” that will occur
%		in Egypt. It will continue to be stated, several times, as a reason for
%		the miraculous plagues.
%	}
}

\Verse{8}
{
	\hP{וְהֵבֵאתִי אֶתְכֶם אֶל הָאָרֶץ אֲשֶׁר נָשָׂאתִי אֶת יָדִי לָתֵת אֹתָהּ לְאַבְרָהָם לְיִצְחָק וּלְיַעֲקֹב וְנָתַתִּי אֹתָהּ לָכֶם מוֹרָשָׁה אֲנִי יְהֹוָה׃}
}
{
	\eP{
		And I shall bring you to the land that I raised my hand to
		give to Abraham, to Isaac, and to Jacob, and I shall give
		it to you as a possession. I am YHWH.’”
	}
}
{
}

\Verse{9}
{
	\hP{וַיְדַבֵּר מֹשֶׁה כֵּן אֶל בְּנֵי יִשְׂרָאֵל וְלֹא שָׁמְעוּ אֶל מֹשֶׁה מִקֹּצֶר רוּחַ וּמֵעֲבֹדָה קָשָׁה׃}
}
{
	\eP{
		And Moses spoke this to the children of Israel, and they
		did not listen to Moses because of shortage of spirit and
		because of hard work.
	}
}
{
}

\Verse{10}
{
	\hP{וַיְדַבֵּר יְהֹוָה אֶל מֹשֶׁה לֵּאמֹר׃}
}
{
	\eP{And YHWH spoke to Moses, saying,}
}
{
}

\Verse{11}
{
	\hP{בֹּא דַבֵּר אֶל פַּרְעֹה מֶלֶךְ מִצְרָיִם וִישַׁלַּח אֶת בְּנֵי יִשְׂרָאֵל מֵאַרְצוֹ׃}
}
{
	\eP{
		“Come, speak to Pharaoh, king of Egypt, that he should let
		the children of Israel go from his land.”
	}
}
{
}

\Verse{12}
{
	\hP{וַיְדַבֵּר מֹשֶׁה לִפְנֵי יְהֹוָה לֵאמֹר הֵן בְּנֵי יִשְׂרָאֵל לֹא שָׁמְעוּ אֵלַי וְאֵיךְ יִשְׁמָעֵנִי פַרְעֹה וַאֲנִי עֲרַל שְׂפָתָיִם׃}
}
{
	\eP{
		And Moses spoke in front of YHWH, saying, “Here, the
		children of Israel didn’t listen to me, and how will
		Pharaoh listen to me?! And I’m uncircumcised of lips!”
		6:12,30. uncircumcised of lips. Does this shed light on
		the meaning of “heavy of tongue and heavy of mouth”? It
		seems rather to be even harder to interpret.
		Uncircumcision is used elsewhere in the Torah as a
		metaphor as well. Lev 26:41 and Deut 10:16 refer to the
		people’s uncircumcised heart. (The prophets Jeremiah and
		Ezekiel use this image also.) There it is a strong image,
		criticizing the people who understand that their covenant
		requires circumcision of the skin that covers the penis,
		but who do not appreciate that they must likewise remove
		the impediment to their hearts. But the meaning of the
		image here is less clear. We still do not know what the
		impediment is. All we can say is that Moses is again
		pictured as unconfident, feeling inadequate to the task.
		The repeated emphasis on this point will make the eventual
		success of the exodus that much more dramatic.
	}
}
{
}

\Verse{13}
{
	\hP{וַיְדַבֵּר יְהֹוָה אֶל מֹשֶׁה וְאֶל אַהֲרֹן וַיְצַוֵּם אֶל בְּנֵי יִשְׂרָאֵל וְאֶל פַּרְעֹה מֶלֶךְ מִצְרָיִם לְהוֹצִיא אֶת בְּנֵי יִשְׂרָאֵל מֵאֶרֶץ מִצְרָיִם׃}
}
{
	\eP{
		And YHWH spoke to Moses and to Aaron and commanded them
		regarding the children of Israel and regarding Pharaoh,
		king of Egypt, to bring out the children of Israel from
		the land of Egypt.
	}
}
{
}

\Verse{14}
{
	\hP{אֵלֶּה רָאשֵׁי בֵית אֲבֹתָם בְּנֵי רְאוּבֵן בְּכֹר יִשְׂרָאֵל חֲנוֹךְ וּפַלּוּא חֶצְרֹן וְכַרְמִי אֵלֶּה מִשְׁפְּחֹת רְאוּבֵן׃}
}
{
	\eP{
		These are the heads of their fathers’ houses: The sons of
		Reuben, Israel’s firstborn: Hanoch and Pallu, Hezron, and
		Carmi. These are the families of Reuben.
	}
}
{
}

\Verse{15}
{
	\hP{וּבְנֵי שִׁמְעוֹן יְמוּאֵל וְיָמִין וְאֹהַד וְיָכִין וְצֹחַר וְשָׁאוּל בֶּן הַכְּנַעֲנִית אֵלֶּה מִשְׁפְּחֹת שִׁמְעוֹן׃}
}
{
	\eP{
		And the sons of Simeon: Jemuel and Jamin and Ohad and
		Jachin and Zohar and Saul, son of a Canaanite woman. These
		are the families of Simeon.
	}
}
{
}

\Verse{16}
{
	\hP{וְאֵלֶּה שְׁמוֹת בְּנֵי לֵוִי לְתֹלְדֹתָם גֵּרְשׁוֹן וּקְהָת וּמְרָרִי וּשְׁנֵי חַיֵּי לֵוִי שֶׁבַע וּשְׁלֹשִׁים וּמְאַת שָׁנָה׃}
}
{
	\eP{
		And these are the names of the sons of Levi by their
		records: Gershon and Kohath and Merari. And the years of
		Levi’s life were a hundred thirty-seven years.
	}
}
{
}

\Verse{17}
{
	\hP{בְּנֵי גֵרְשׁוֹן לִבְנִי וְשִׁמְעִי לְמִשְׁפְּחֹתָם׃}
}
{
	\eP{The sons of Gershon: Libni and Shimei, by their families.}
}
{
}

\Verse{18}
{
	\hP{וּבְנֵי קְהָת עַמְרָם וְיִצְהָר וְחֶבְרוֹן וְעֻזִּיאֵל וּשְׁנֵי חַיֵּי קְהָת שָׁלֹשׁ וּשְׁלֹשִׁים וּמְאַת שָׁנָה׃}
}
{
	\eP{
		And the sons of Kohath: Amram and Izhar and Hebron and
		Uzziel. And the years of Kohath’s life were a hundred
		thirty-three years.
	}
}
{
}

\Verse{19}
{
	\hP{וּבְנֵי מְרָרִי מַחְלִי וּמוּשִׁי אֵלֶּה מִשְׁפְּחֹת הַלֵּוִי לְתֹלְדֹתָם׃}
}
{
	\eP{
		And the sons of Merari: Mahli and Mushi. These are the
		families of Levi by their records.
	}
}
{
}

\Verse{20}
{
	\hP{וַיִּקַּח עַמְרָם אֶת יוֹכֶבֶד דֹּדָתוֹ לוֹ לְאִשָּׁה וַתֵּלֶד לוֹ אֶת אַהֲרֹן וְאֶת מֹשֶׁה וּשְׁנֵי חַיֵּי עַמְרָם שֶׁבַע וּשְׁלֹשִׁים וּמְאַת שָׁנָה׃}
}
{
	\eP{
		And Amram took Jochebed, his aunt, as a wife; and she gave
		birth to Aaron and Moses for him. And the years of Amram’s
		life were a hundred thirty-seven years.
	}
}
{
%	\footnote{
%		she gave birth to Aaron and Moses. Until now, it has not been clear that
%		Moses and Aaron are brothers. There is no reference to a brother in the
%		story of Moses’ birth, and at the bush Aaron is referred to as “your
%		Levite brother,” which sounds like a fellow Levite, not an actual
%		brother. The list in which this appears (6:13–25) seems strange: it
%		comes in the middle of a conversation between God and Moses, and it
%		covers only Reuben, Simeon, and Levi, and then leaves out the rest of
%		the tribes. But it appears that the point of the list was to get to
%		Aaron: to identify him as Moses’ brother, to state that he is married to
%		the sister of the chieftain of the tribe of Judah, and to name his son
%		and grandson who will be his successors. The order also indicates that
%		Aaron is the firstborn; this is confirmed in Exod 7:7. Footnote 6:20.
%		Aaron and Moses. And what about Miriam? One Hebrew manuscript and the
%		Septuagint and Samaritan texts have the words “and their sister Miriam.”
%		The scribe appears to have accidentally omitted these words here.
%	}
}

\Verse{21}
{
	\hP{וּבְנֵי יִצְהָר קֹרַח וָנֶפֶג וְזִכְרִי׃}
}
{
	\eP{And the sons of Izhar: Korah and Nepheg and Zichri.}
}
{
}

\Verse{22}
{
	\hP{וּבְנֵי עֻזִּיאֵל מִישָׁאֵל וְאֶלְצָפָן וְסִתְרִי׃}
}
{
	\eP{And the sons of Uzziel: Mishael and Elzaphan and Sithri.}
}
{
}

\Verse{23}
{
	\hP{וַיִּקַּח אַהֲרֹן אֶת אֱלִישֶׁבַע בַּת עַמִּינָדָב אֲחוֹת נַחְשׁוֹן לוֹ לְאִשָּׁה וַתֵּלֶד לוֹ אֶת נָדָב וְאֶת אֲבִיהוּא אֶת אֶלְעָזָר וְאֶת אִיתָמָר׃}
}
{
	\eP{
		And Aaron took Elisheba, daughter of Amminadab, sister of
		Nahshon as a wife; and she gave birth to Nadab and Abihu,
		Eleazar, and Ithamar for him.
	}
}
{
%	\footnote{
%		daughter of Amminadab, sister of Nahshon. Aaron marries the sister of
%		the chieftain of the tribe of Judah. It is noteworthy that the priestly
%		and political elite marry (and this suggests that such marriages occured
%		later between the royal and high priestly families of Israel).
%	}
}

\Verse{24}
{
	\hP{וּבְנֵי קֹרַח אַסִּיר וְאֶלְקָנָה וַאֲבִיאָסָף אֵלֶּה מִשְׁפְּחֹת הַקּרְחִי׃}
}
{
	\eP{
		And the sons of Korah: Assir and Elkanah and Abiasaph.
		These are the families of the Korahites.
	}
}
{
}

\Verse{25}
{
	\hP{וְאֶלְעָזָר בֶּן אַהֲרֹן לָקַח לוֹ מִבְּנוֹת פּוּטִיאֵל לוֹ לְאִשָּׁה וַתֵּלֶד לוֹ אֶת פִּינְחָס אֵלֶּה רָאשֵׁי אֲבוֹת הַלְוִיִּם לְמִשְׁפְּחֹתָם׃}
}
{
	\eP{
		And Eleazar son of Aaron had taken one of the daughters of
		Putiel as a wife, and she gave birth for him to Phinehas.
		These are the heads of the fathers of the Levites
		according to their families.
	}
}
{
}

\Verse{26}
{
	\hP{הוּא אַהֲרֹן וּמֹשֶׁה אֲשֶׁר אָמַר יְהֹוָה לָהֶם הוֹצִיאוּ אֶת בְּנֵי יִשְׂרָאֵל מֵאֶרֶץ מִצְרַיִם עַל צִבְאֹתָם׃}
}
{
	\eP{
		That is Aaron and Moses, to whom YHWH said, “Bring out the
		children of Israel from the land of Egypt by their
		masses.” 6:26–27. Aaron and Moses … Moses and Aaron. Both
		word orders are given: Aaron first, then Moses first,
		suggesting that each is important in a different way. The
		recognition of Aaron here is especially significant
		because it is natural to think of Moses first, but this
		notice follows a genealogy that has listed Aaron as the
		firstborn and identified his wife and children. (The
		Septuagint and Qumran texts have the order in v. 27
		reversed, which unfortunately loses this point.)
	}
}
{
}

\Verse{27}
{
	\hP{הֵם הַמְדַבְּרִים אֶל פַּרְעֹה מֶלֶךְ מִצְרַיִם לְהוֹצִיא אֶת בְּנֵי יִשְׂרָאֵל מִמִּצְרָיִם הוּא מֹשֶׁה וְאַהֲרֹן׃}
}
{
	\eP{
		They were the ones speaking to Pharaoh, king of Egypt, to
		bring out the children of Israel from Egypt. That is Moses
		and Aaron.
	}
}
{
}

\Verse{28}
{
	\hP{וַיְהִי בְּיוֹם דִּבֶּר יְהֹוָה אֶל מֹשֶׁה בְּאֶרֶץ מִצְרָיִם׃}
}
{
	\eP{
		And it was in the day that YHWH spoke to Moses in the land
		of Egypt,
	}
}
{
}

\Verse{29}
{
	\hP{וַיְדַבֵּר יְהֹוָה אֶל מֹשֶׁה לֵּאמֹר אֲנִי יְהֹוָה דַּבֵּר אֶל פַּרְעֹה מֶלֶךְ מִצְרַיִם אֵת כּל אֲשֶׁר אֲנִי דֹּבֵר אֵלֶיךָ׃}
}
{
	\eP{
		and YHWH spoke to Moses, saying, “I am YHWH. Speak to
		Pharaoh, king of Egypt, everything that I speak to you.”
	}
}
{
}

\Verse{30}
{
	\hP{וַיֹּאמֶר מֹשֶׁה לִפְנֵי יְהֹוָה הֵן אֲנִי עֲרַל שְׂפָתַיִם וְאֵיךְ יִשְׁמַע אֵלַי פַּרְעֹה׃}
}
{
	\eP{
		And Moses said in front of YHWH, “Here, I’m uncircumcised
		of lips, and how will Pharaoh listen to me?!”
	}
}
{
}
